\documentclass[notes=show]{beamer}
\usepackage{mathpazo}
\usepackage{hyperref}
\usepackage{multimedia}
\usepackage{tcolorbox}
\usepackage{tikz}
\usetikzlibrary{shadows}
\usetheme{metropolis}
\usecolortheme{default}
\setbeamertemplate{itemize item}{%
    \begin{tikzpicture}
        \shade[ball color=black!100!yellow, preaction={fill=black,
opacity=.00,transform canvas={xshift=1mm,yshift=-1mm, yscale=0.5}}] (0,0) circle (0.6ex);
    \end{tikzpicture}
}
\setbeamertemplate{itemize subitem}{%
    \begin{tikzpicture}
        \shade[ball color=black!100!white, preaction={fill=black,
        opacity=.00,transform canvas={xshift=1mm,yshift=-1mm, yscale=0.5}}] (0,0) circle (0.6ex);
    \end{tikzpicture}
}
\setbeamercolor{frametitle}{bg=white}
\setbeamercolor{frametitle}{fg=black}
\setbeamercolor{background canvas}{bg=white}
\setbeamercolor{block body}{bg=mDarkTeal!30}
\setbeamercolor{block title}{bg=mDarkTeal,fg=black!2}
\setbeamertemplate{navigation symbols}{}
\setbeamertemplate{footline}[page number]
\newenvironment{stepenumerate}{\begin{enumerate}[<+->]}{\end{enumerate}}
\newenvironment{stepitemize}{\begin{itemize}[<+->]}{\end{itemize} }
\newenvironment{stepenumeratewithalert}{\begin{enumerate}[<+-| alert@+>]}{\end{enumerate}}
\newenvironment{stepitemizewithalert}{\begin{itemize}[<+-| alert@+>]}{\end{itemize} }

\graphicspath{{figures/}}

\AtBeginSection[]
{
\begin{frame}{Outline}
    \tableofcontents[currentsection]
\end{frame}
}

\begin{document}

\title{The China Syndrome: Local Labor Market Effects of Import Competition in the US}
\subtitle{}
\date{David Autor, David Dorn \& Gordon Hanson \bigskip \\
\textit{American Economic Review}, October 2013 }
\maketitle

%%%%%%%%%%%%%%%%%%%%%%%%%%%%%%%%%%%%%%%%%
\section{Introduction}
%%%%%%%%%%%%%%%%%%%%%%%%%%%%%%%%%%%%%%%%%

\begin{frame}{Stylized facts}
\begin{itemize}
\item Pre-1990s limited impact trade on US labor \bigskip
\item Trend decline in US manufacturing employment \bigskip
\item Increase in import competition from China \bigskip
\item Varies across regions in the US \bigskip
\item Limited offsetting increase in demand for US exports
\end{itemize}
\end{frame}

\begin{frame}
\frametitle{Import penetration and manufacturing employment}
\begin{center}
\includegraphics[width=\textwidth]{fig-1.pdf}%
\end{center}
\end{frame}

\begin{frame}
\frametitle{Long-run decline in US manufacturing employment}
\begin{center}
\includegraphics[width=\textwidth]{AR-fig-1.pdf}%
\end{center}
\end{frame}

\begin{frame}{Research question}
\begin{itemize}
\item<1-> Impact of Chinese imports on US local labor markets (CZ)
\begin{itemize} \medskip
\item<1-> CZs differently exposed because of industry specialization 
\end{itemize} \bigskip
\item<1-> Analyze impact on non-wage variables 
\begin{itemize} \medskip
\item<1-> Employment, unemployment, participation, mobility, transfers 
\end{itemize} \bigskip
\item<1-> Develop Shift-Share Instrumental Variable (SSIV)
\begin{itemize} \medskip
\item<1-> Instrument US import competition (share $\times$ shock) with imports (shock) by other countries
\end{itemize}
\end{itemize}
\end{frame}

\begin{frame}
\frametitle{Key findings for trade-exposed CZs}
\begin{itemize}
\item<1-> Lower wages and employment in (non-)manufacturing \bigskip
\item<1-> Higher unemployment and lower labor force participation \bigskip
\item<1-> Increased reliance on welfare benefits \bigskip
\item<1-> Limited evidence of emigration \bigskip
\item<1-> Significant distributional consequences of trade in reality!
\end{itemize}
\end{frame}

%%%%%%%%%%%%%%%%%%%%%%%%%%%%%%%%%%%%%%%%%
\section{Theoretical predictions}
%%%%%%%%%%%%%%%%%%%%%%%%%%%%%%%%%%%%%%%%%

\begin{frame}{Exogenous shocks}
\begin{itemize}
\item<1-> China experienced: \bigskip
\begin{itemize}
    \item<1-> Productivity growth (central planning $\rightarrow$ market economy) \bigskip
    \item<1-> Reduction in trade costs (accession to WTO in 2000) \bigskip
\end{itemize} \medskip
\item<1-> Opposing shocks in demand for goods in US CZs: \bigskip
\begin{itemize}
    \item<1-> Increased competition in markets in which CZs sell output \bigskip
    \item<1-> Increased demand for goods in China
\end{itemize}
\end{itemize}
\end{frame}

\begin{frame}
\frametitle{China's rise in world manufacturing}
\begin{center}
\includegraphics[width=\textwidth]{AR-fig-2.pdf}%
\end{center}
\end{frame}

\begin{frame}
\frametitle{China's revealed comparative advantage in manufacturing}
\begin{center}
\includegraphics[width=\textwidth]{AR-fig-3.pdf}%
\end{center}
\end{frame}

\begin{frame}{Impacts on labor demand and supply in US CZs}
\begin{itemize}
\item<1-> Positive shock to China’s export supply: \medskip
\begin{itemize}
    \item<1-> Decrease labor demand in traded goods \medskip
    \item<1-> Increase labor supply and demand in non-traded sectors \medskip
\end{itemize}
\item<1-> Positive shock to China’s import demand: \medskip
\begin{itemize}
    \item<1-> Increase labor demand in traded goods \medskip
    \item<1-> Decrease labor supply and increase labor demand in non-traded sectors \medskip
\end{itemize}
\item<1-> Overall adverse effects driven by US trade deficit and regional specialization within US
\end{itemize}
\end{frame}


\begin{frame}
\frametitle{Impact on manufacturers of imports}
\begin{center}
\includegraphics[width=\textwidth]{fig-model-1.pdf}%
\end{center}
\end{frame}

\begin{frame}
\frametitle{Impact on manufacturers of exports}
\begin{center}
\includegraphics[width=\textwidth]{fig-model-2.pdf}%
\end{center}
\end{frame}

\begin{frame}
\frametitle{Current account balances}
\begin{center}
\includegraphics[width=\textwidth]{AR-fig-5.pdf}%
\end{center}
\end{frame}

%%%%%%%%%%%%%%%%%%%%%%%%%%%%%%%%%%%%%%%%%
\section{Measuring exposure to import competition}
%%%%%%%%%%%%%%%%%%%%%%%%%%%%%%%%%%%%%%%%%

\begin{frame}{Measuring exposure to import competition}
\begin{itemize}
\item Change in Chinese import exposure per worker in a CZ:
\begin{equation*}
    \Delta IPW_{zt} = \frac{\sum_{i} L_{izt} \Delta m_{it}}{L_{zt}} \text{ with } \Delta m_{it} = \frac{\Delta M_{it}}{L_{it}}
\end{equation*}
for CZ $z$ and industry $i$ \medskip
\item Rewrite this as:
\begin{equation*}
    \Delta IPW_{zt} = \sum_{i} \frac{L_{izt}}{L_{zt}} \frac{\Delta M_{it}}{L_{it}} \tag{3}
\end{equation*}
with $L_{izt} / L_{zt}$ the degree of industrial specialization in $z$ and $\Delta M_{it} / L_{it}$ the import value per worker in industry $i$
\end{itemize}
\end{frame}

\begin{frame}
\frametitle{Trade exposure across CZs}
\begin{center}
\includegraphics[width=\textwidth]{AR-fig-6.pdf}%
\end{center}
\end{frame}

\begin{frame}
\frametitle{Change in net import penetration by sector}
\begin{center}
\includegraphics[width=\textwidth]{AR-fig-4.pdf}%
\end{center}
\end{frame}

\begin{frame}{Shift-share instrument for import competition}
\begin{itemize}
\item $\Delta m_{it} $ could be correlated with unobserved US demand shocks \medskip
\item Define the instrument:
\begin{equation*}
    \Delta \tilde{IPW}_{zt} = \sum_{i} \frac{L_{izt-1}}{L_{zt-1}} \frac{\Delta \tilde{M}_{it}}{L_{it-1}} \tag{4}
\end{equation*}
with $ \Delta \tilde{M}_{it}$ Chinese imports into 8 other countries and $t-1$ to account for expected employment impact of China shock at $t$ \medskip
\item Instrument predicts supply-driven changes in Chinese imports (e.g. China's SEZ, accession to WTO in 2000)
\end{itemize}
\end{frame}

\begin{frame}{Empirical research design}
\begin{itemize}
\item SSIV using a 2SLS estimator (using an instrument for the shift $\Delta m_{it} $ and not for the share $L_{izt} / L_{zt}$)  \medskip
\item Second stage: 
\begin{equation*}
  \Delta Y_{zt}^m = \gamma_{t} + \beta_{1} \Delta \hat{IPW}_{zt} + X_{it}' \beta_{2} + \epsilon_{zt} \tag{5}
\end{equation*}
\item First stage:
\begin{equation*}
  \Delta IPW_{zt} = \beta_{3} \Delta \tilde{IPW}_{zt} + \tilde{\epsilon}_{zt}
\end{equation*}
\item Standard errors are clustered at the state level to allow for spatial correlations across CZs
\end{itemize}
\end{frame}

%%%%%%%%%%%%%%%%%%%%%%%%%%%%%%%%%%%%%%%%%
\section{Results}
%%%%%%%%%%%%%%%%%%%%%%%%%%%%%%%%%%%%%%%%%

%%%%%%%%%%%%%%%%%%%%%%%%%%%%%%%%%%%%%%%%%
\subsection{Adverse impact on manufacturing employment}
%%%%%%%%%%%%%%%%%%%%%%%%%%%%%%%%%%%%%%%%%

\begin{frame}
\frametitle{First-stage estimates}
\begin{center}
\includegraphics[width=\textwidth]{fig-2A.pdf}%
\end{center}
\end{frame}

\begin{frame}
\frametitle{Reduced-form estimates}
\begin{center}
\includegraphics[width=\textwidth]{fig-2B.pdf}%
\end{center}
\end{frame}

\begin{frame}
\frametitle{Decline in manufacturing employment (no controls)}
\begin{center}
\includegraphics[width=\textwidth]{table-2.pdf}%
\end{center}
\end{frame}

\begin{frame}
\frametitle{Decline in manufacturing employment (with controls)}
\begin{center}
\includegraphics[width=\textwidth]{table-3.pdf}%
\end{center}
\end{frame}

%%%%%%%%%%%%%%%%%%%%%%%%%%%%%%%%%%%%%%%%%
\subsection{Limited worker reallocation between CZs}
%%%%%%%%%%%%%%%%%%%%%%%%%%%%%%%%%%%%%%%%%

\begin{frame}
\frametitle{Limited worker reallocation between CZs}
\begin{center}
\includegraphics[width=\textwidth]{table-4.pdf}%
\end{center}
\end{frame}

\begin{frame}{Limited reallocation of workers between CZ}
\begin{itemize}
\item No evidence that shocks to local manufacturing lead to change in population \bigskip
\item Shocks to manufacturing too small to affect outcomes in broader CZ \bigskip
\item Shocks are substantial but no reallocation of workers because moving costs too high \bigskip
\item CZ absorbs shock through other employment in CZ and/or lower wages 
\end{itemize}
\end{frame}

%%%%%%%%%%%%%%%%%%%%%%%%%%%%%%%%%%%%%%%%%
\subsection{Adjustments within trade-exposed CZs}
%%%%%%%%%%%%%%%%%%%%%%%%%%%%%%%%%%%%%%%%%

\begin{frame}
\frametitle{Job loss in trade-exposed CZs}
\begin{center}
\includegraphics[width=\textwidth]{table-5.pdf}%
\end{center}
\end{frame}

\begin{frame}
\frametitle{Negative impact on average wage in trade-exposed CZs}
\begin{center}
\includegraphics[width=\textwidth]{table-6.pdf}%
\end{center}
\end{frame}

\begin{frame}
\frametitle{Negative impact on non-manuf. wage in trade-exposed CZs}
\begin{center}
\includegraphics[width=\textwidth]{table-7.pdf}%
\end{center}
\end{frame}

\begin{frame}
\frametitle{Increased welfare payments in trade-exposed CZs}
\begin{center}
\includegraphics[width=\textwidth]{table-8.pdf}%
\end{center}
\end{frame}


\begin{frame}
\frametitle{Impact of \$1000 increase in exposure on welfare payments}
\begin{center}
\includegraphics[width=\textwidth]{AR-fig-7.pdf}%
\end{center}
\end{frame}

\begin{frame}
\frametitle{Decrease in HH income in trade-exposed CZs}
\begin{center}
\includegraphics[width=\textwidth]{table-9.pdf}%
\end{center}
\end{frame}

%%%%%%%%%%%%%%%%%%%%%%%%%%%%%%%%%%%%%%%%%
\section{Robustness checks}
%%%%%%%%%%%%%%%%%%%%%%%%%%%%%%%%%%%%%%%%%

\begin{frame}{Different measures of import competition}
\begin{itemize}
    \item Include competition in foreign markets \bigskip
    \item Include imports of intermediate inputs \bigskip
    \item Use net import from China (i.e. include export effects) \bigskip
    \item Use gravity residual of relative productivity changes \bigskip
    \item Use H-O factor content of US net imports from China
\end{itemize}
\end{frame}

%%%%%%%%%%%%%%%%%%%%%%%%%%%%%%%%%%%%%%%%%
\section*{Conclusions}
%%%%%%%%%%%%%%%%%%%%%%%%%%%%%%%%%%%%%%%%%

\begin{frame}{Conclusions}
\begin{itemize}
\item Impact of Chinese import competition through reducing manufacturing employment \bigskip
\item Chinese import competition explains 21 percent of decline manufacturing employment in U.S. from 1991 to 2007 \bigskip
\item Limited labor market adjustments across CZ \bigskip
\item In trade-exposed CZ, there is a decline in non-manufacturing wages, increase in unemployment and transfer payments
\end{itemize}
\end{frame}


%%%%%%%%%%%%%%%%%%%%%%%%%%%%%%%%%%%%%%%%%
\section*{Further readings}
%%%%%%%%%%%%%%%%%%%%%%%%%%%%%%%%%%%%%%%%%

\begin{frame}{Further readings}
\begin{itemize}
\footnotesize
\item Autor, Dorn, Hanson, and Song. 2014. Trade adjustment: Worker Level Evidence. \textit{Quartely Journal of Economics}. 129(4):1799–860 \smallskip
\item Autor, Dorn, and Hanson. 2016. The China Shock: Learning from Labor-Market Adjustment to Large Changes in Trade. \textit{Annual Review of Economics}. 8:205–40 \medskip
\item Autor, Dorn, and Hanson. 2019. When Work Disappears: Manufacturing Decline and the Falling Marriage Market Value of Young Men. \textit{American Economic Review: Insights}. 1(2):161-178 \medskip
\item Autor, Dorn, Hanson, and Majlesi. 2020. Importing Political Polarization? The Electoral Consequences of Rising Trade Exposure. \textit{American Economic Review}. 110(10):3139–3183 \medskip
\item Autor, Dorn, and Hanson. 2021. On the Persistence of the China Shock. forthcoming in \textit{Brookings Papers on Economic Activity}. \medskip
\item Autor, Beck, Dorn, and Hanson. 2024. Help for the Heartland? The Employment and Electoral Effects of the Trump Tariffs.
\end{itemize}
\end{frame}

\end{document}
